\chapter{Sürekli Fonksiyonlar ve Homeomorfizmalar}

Topolojik uzaylar arasındaki fonksiyonlarda süreklilik, açık kümelerin
geri çekiminin açık olması koşuluna dayanır.

%-------------------------------------------------------------------
\section{Konu}

\begin{sezgi}
Sürekliliği ``açıklığın geri taşınması'' gibi düşünün: hedef uzayda hangi
açık kümeye bakarsanız bakın, onun \emph{geri çekimi} (preimage) kaynak
uzayda yine açık olmalı. Analizden bildiğiniz $\varepsilon$--$\delta$ tanımı
bunun metrik özelidir; topolojide $\varepsilon$--$\delta$ yerine doğrudan
açık kümelerin dilini kullanırız. Homeomorfizma ise bu taşımanın \emph{iki
yönde} de bozulmadan işlemesi --- yani iki uzayın topolojik olarak ``aynı''
olmasıdır.
\end{sezgi}

\begin{definition}[Sürekli Fonksiyon]
$f: (X, \tau_X) \to (Y, \tau_Y)$ fonksiyonu \textbf{sürekli} ise:
\[
\forall V \in \tau_Y:\; f^{-1}(V) \in \tau_X.
\]
\end{definition}

\begin{figure}[ht]
  \centering
  \begin{tikzpicture}[font=\small]
  % --- domain X ---
  \draw[rounded corners=8pt, thick] (-0.4,-1.3) rectangle (3.0,1.5);
  \node at (-0.05,1.2) {$X$};
  \draw[blue!70!black, thick, fill=blue!10] (1.3,-0.1) ellipse [x radius=0.95, y radius=0.7];
  \node[blue!70!black] at (1.3,-0.1) {$f^{-1}(V)$};
  % --- codomain Y ---
  \draw[rounded corners=8pt, thick] (5.0,-1.3) rectangle (8.4,1.5);
  \node at (8.05,1.2) {$Y$};
  \draw[violet!80!black, thick, fill=violet!10] (6.7,-0.1) ellipse [x radius=0.85, y radius=0.6];
  \node[violet!80!black] at (6.7,-0.1) {$V$};
  % --- arrow f ---
  \draw[-{Stealth}, thick] (3.25,0.55) to[bend left=12]
    node[above, pos=0.5] {$f$} (4.75,0.55);
  % --- caption note ---
  \node[font=\scriptsize, align=center] at (4.0,-1.85)
    {$V$ açık $\Rightarrow f^{-1}(V)$ açık olmalı (her açık $V$ için)};
\end{tikzpicture}

  \caption{Süreklilik: hedef uzaydaki her açık $V$ için geri çekim
    $f^{-1}(V)$ kaynak uzayda açık olmalıdır.}
  \label{fig:ch10:sureklilik}
\end{figure}

Eşdeğer koşullar:
\begin{itemize}
  \item Her kapalı $F \subseteq Y$ için $f^{-1}(F)$ kapalıdır.
  \item Her $x \in X$ ve $f(x)$'in her komşuluğu $W$ için $f^{-1}(W)$, $x$'in komşuluğudur.
\end{itemize}

\begin{definition}[Homeomorfizma]
$f: X \to Y$ \textbf{homeomorfizma} ise: bijektif, $f$ sürekli, $f^{-1}$ sürekli.
$(X, \tau_X) \cong (Y, \tau_Y)$ ile gösterilir.
\end{definition}

\begin{figure}[ht]
  \centering
  \begin{tikzpicture}[font=\small]
  % --- space X ---
  \draw[blue!70!black, thick, fill=blue!8] (0,0) ellipse [x radius=1.2, y radius=0.95];
  \node[blue!70!black] at (0,1.3) {$(X,\tau_X)$};
  % --- space Y ---
  \draw[violet!80!black, thick, fill=violet!8] (5.0,0) ellipse [x radius=1.2, y radius=0.95];
  \node[violet!80!black] at (5.0,1.3) {$(Y,\tau_Y)$};
  % --- two arrows ---
  \draw[-{Stealth}, thick] (1.3,0.35) to[bend left=22]
    node[above, pos=0.5] {$f$ sürekli} (3.7,0.35);
  \draw[-{Stealth}, thick] (3.7,-0.35) to[bend left=22]
    node[below, pos=0.5] {$f^{-1}$ sürekli} (1.3,-0.35);
  % --- caption ---
  \node[font=\scriptsize, align=center] at (2.5,-1.7)
    {bijektif, iki yön de sürekli $\Rightarrow (X,\tau_X)\cong(Y,\tau_Y)$};
\end{tikzpicture}

  \caption{Homeomorfizma: bijektif bir eşleme hem $f$ hem $f^{-1}$ yönünde
    sürekliyse iki uzay topolojik olarak özdeştir, $(X,\tau_X)\cong(Y,\tau_Y)$.}
  \label{fig:ch10:homeo}
\end{figure}

\begin{karsiornek}
Süreklilik \emph{yöne duyarlıdır}. Aynı taşıyıcı küme $\{0,1\}$ üzerinde
özdeşlik fonksiyonunu düşünün: ayrık topolojiden (ince) Sierpiński'ye (kaba)
\textbf{sürekli}, ama Sierpiński'den ayrığa \textbf{sürekli değil} --- çünkü
ayrıktaki açık $\{0\}$'ın geri çekimi Sierpiński'de açık değildir. Yani
``kaba $\to$ ince'' yönünde özdeşlik genelde süreklilik kaybeder
(Örnek~5.8'de \texttt{is\_continuous\_finite\_map} ile doğrulanır).
\end{karsiornek}

\begin{table}[h]
\centering
\begin{tabular}{ll}
\toprule
\textbf{Tür} & \textbf{Koşul} \\
\midrule
Sabit fonksiyon & $f(x) = c$; her zaman sürekli \\
Özdeşlik & $f(x) = x$; her zaman sürekli \\
Bileşke & $f$, $g$ sürekli $\Rightarrow g \circ f$ sürekli \\
Homeomorfizma & Bijektif; $f$ ve $f^{-1}$ sürekli \\
\bottomrule
\end{tabular}
\caption{Süreklilik hiyerarşisi}
\end{table}

%-------------------------------------------------------------------
\section{Teoremler}

\begin{theorem}
Her sabit fonksiyon süreklidir.
\begin{proof}[Kanıt İskeleti]
$f^{-1}(V) = X$ eğer $c \in V$, $\emptyset$ eğer $c \notin V$; her ikisi de açıktır.
\end{proof}
\end{theorem}

\begin{theorem}
$f: X \to Y$ ve $g: Y \to Z$ sürekli $\Rightarrow g \circ f$ süreklidir.
\begin{proof}[İspat eskizi]
$W \subseteq Z$ açık olsun. $g$ sürekli olduğundan $g^{-1}(W) \subseteq Y$
açıktır. $f$ sürekli olduğundan $f^{-1}(g^{-1}(W)) \subseteq X$ açıktır.
$(g\circ f)^{-1}(W) = f^{-1}(g^{-1}(W))$ olduğundan $g\circ f$'in her açık
kümenin geri çekimi açıktır; yani $g\circ f$ süreklidir.
\end{proof}
\end{theorem}

\begin{figure}[ht]
  \centering
  \begin{tikzpicture}[font=\small]
  \draw[blue!70!black, thick, fill=blue!8] (0,0) ellipse [x radius=0.95, y radius=0.85];
  \node[blue!70!black] at (0,1.15) {$X$};
  \draw[teal!70!black, thick, fill=teal!8] (3.5,0) ellipse [x radius=0.95, y radius=0.85];
  \node[teal!70!black] at (3.5,1.15) {$Y$};
  \draw[violet!80!black, thick, fill=violet!8] (7.0,0) ellipse [x radius=0.95, y radius=0.85];
  \node[violet!80!black] at (7.0,1.15) {$Z$};
  \draw[-{Stealth}, thick] (1.05,0.0) -- node[above] {$f$} (2.45,0.0);
  \draw[-{Stealth}, thick] (4.55,0.0) -- node[above] {$g$} (5.95,0.0);
  \draw[-{Stealth}, thick, red!70!black] (0.3,-0.75) to[bend right=24]
    node[below, pos=0.5] {$g\circ f$} (6.7,-0.75);
  \node[font=\scriptsize, align=center] at (3.5,-2.0)
    {$f$ ve $g$ sürekli $\Rightarrow g\circ f$ sürekli};
\end{tikzpicture}

  \caption{Bileşke süreklilik: $f$ ve $g$ sürekliyse $g\circ f$ de süreklidir;
    geri çekim zincirlenir, $(g\circ f)^{-1}(W) = f^{-1}(g^{-1}(W))$.}
  \label{fig:ch10:bileske}
\end{figure}

\begin{theorem}
$f: X \to Y$ sürekli, $X$ kompakt $\Rightarrow f(X)$ kompakttır.
\begin{proof}[İspat eskizi]
$f(X)$'in bir açık örtüsü $\{V_\alpha\}$ verilsin. $f$ sürekli olduğundan her
$f^{-1}(V_\alpha)$ $X$'te açıktır ve bunlar $X$'i örter. $X$ kompakt olduğundan
sonlu bir alt-örtü $f^{-1}(V_{\alpha_1}), \dots, f^{-1}(V_{\alpha_n})$ vardır.
O zaman $V_{\alpha_1}, \dots, V_{\alpha_n}$ $f(X)$'i örter: süreklilik açık
örtüyü geri çeker, kompaktlık sonlu alt-örtü verir, görüntüye geri itilir.
Demek ki $f(X)$ kompakttır.
\end{proof}
\end{theorem}

\begin{theorem}
$f: X \to Y$ sürekli, $X$ bağlantılı $\Rightarrow f(X)$ bağlantılıdır.
\end{theorem}

\begin{theorem}
Kompakt $X$'ten Hausdorff $Y$'ye sürekli bijeksiyon $\Rightarrow$ homeomorfizma.
\end{theorem}

%-------------------------------------------------------------------
\section{Algoritmalar}

\subsection{is\_continuous\_finite\_map — $O(|\tau_Y| \cdot |X|)$}

\begin{algorithm}
\caption{IsContinuous$(X, \tau_X, Y, \tau_Y, f)$}
\begin{algorithmic}[1]
\For{each $V \in \tau_Y$}
    \State preimage $\leftarrow \{x \in X : f(x) \in V\}$
    \If{preimage $\notin \tau_X$}
        \State \Return False
    \EndIf
\EndFor
\State \Return True
\end{algorithmic}
\end{algorithm}

%-------------------------------------------------------------------
\section{pytop API}

\begin{lstlisting}[language=Python]
from pytop import (
    is_continuous_finite_map,
    finite_homeomorphism_result,
    sierpinski_space,
    discrete_topology,
    indiscrete_topology,
    make_topology,
    make_set,
    empty_set,
)
\end{lstlisting}

\texttt{is\_continuous\_finite\_map(domain, domain\_topology, codomain, codomain\_topology, mapping)} $\to$ \texttt{bool}

\texttt{finite\_homeomorphism\_result(left, right)} $\to$ \texttt{Result}

%-------------------------------------------------------------------
\section{Örnekler}

\subsection*{Örnek 5.2 — Ayrık Topoloji: Her Fonksiyon Sürekli}

\begin{lstlisting}[language=Python]
d = discrete_topology(0, 1)
d_pts = list(d.carrier)
d_topo = [set(u) for u in d.topology]

f_id   = {0: 0, 1: 1}
f_swap = {0: 1, 1: 0}
print("ozdeslik:", is_continuous_finite_map(d_pts, d_topo, d_pts, d_topo, f_id))
print("swap:    ", is_continuous_finite_map(d_pts, d_topo, d_pts, d_topo, f_swap))
\end{lstlisting}

\begin{verbatim}
ozdeslik continuous: True
swap     continuous: True
\end{verbatim}

\subsection*{Örnek 5.3 — Sierpiński $\to$ Ayrık: Sürekli Değil}

\begin{lstlisting}[language=Python]
s = sierpinski_space()
s_pts = list(s.carrier)
s_topo = [set(u) for u in s.topology]
print("id: S->D:", is_continuous_finite_map(s_pts, s_topo, d_pts, d_topo, f_id))
print("id: D->S:", is_continuous_finite_map(d_pts, d_topo, s_pts, s_topo, f_id))
\end{lstlisting}

\begin{verbatim}
id: Sierpinski -> Disc continuous: False
id: Disc -> Sierpinski continuous: True
\end{verbatim}

Sierpiński $\tau = \{\emptyset, \{0,1\}, \{1\}\}$: $f^{-1}(\{0\}) = \{0\}$ Sierpiński'de açık değil.

\subsection*{Örnek 5.5 — Homeomorfizma Kontrolü}

\begin{lstlisting}[language=Python]
d2a = discrete_topology(1, 2)
d2b = discrete_topology('a', 'b')
ind2 = indiscrete_topology(1, 2)
print("D(1,2) ~ D(a,b):", finite_homeomorphism_result(d2a, d2b).status)
print("D(1,2) ~ S(0,1):", finite_homeomorphism_result(d2a, s).status)
print("D(1,2) ~ Ind(1,2):", finite_homeomorphism_result(d2a, ind2).status)
\end{lstlisting}

\begin{verbatim}
D(1,2) ~ D(a,b): true
D(1,2) ~ S(0,1): false
D(1,2) ~ Ind(1,2): false
\end{verbatim}

\subsection*{Örnek 5.7 — Bileşke Süreklilik: $g \circ f$}

\begin{lstlisting}[language=Python]
TX = make_topology(make_set(1, 2, 3), make_set(1), make_set(1, 2))
TY = make_topology(make_set('a', 'b', 'c'), make_set('a'), make_set('a', 'b'))
TZ = make_topology(make_set('p', 'q', 'r'), make_set('p'), make_set('p', 'q'))
X_pts, X_topo = list(TX.carrier), list(TX.topology)
Y_pts, Y_topo = list(TY.carrier), list(TY.topology)
Z_pts, Z_topo = list(TZ.carrier), list(TZ.topology)

f = {1: 'a', 2: 'b', 3: 'c'}
g = {'a': 'p', 'b': 'q', 'c': 'r'}
gof = {x: g[f[x]] for x in X_pts}
print("f continuous:    ", is_continuous_finite_map(X_pts, X_topo, Y_pts, Y_topo, f))
print("g continuous:    ", is_continuous_finite_map(Y_pts, Y_topo, Z_pts, Z_topo, g))
print("g of f continuous:", is_continuous_finite_map(X_pts, X_topo, Z_pts, Z_topo, gof))
\end{lstlisting}

\begin{verbatim}
f continuous:     True
g continuous:     True
g of f continuous: True
\end{verbatim}

İki sürekli fonksiyonun geri çekimleri zincirlenir:
$(g \circ f)^{-1}(W) = f^{-1}(g^{-1}(W))$ (Teorem~2.2'nin sayısal doğrulaması).

\subsection*{Örnek 5.8 — Süreklilik Yöne Duyarlı}

\begin{lstlisting}[language=Python]
fine = discrete_topology(0, 1)
coarse = sierpinski_space()
fine_pts, fine_topo = list(fine.carrier), list(fine.topology)
co_pts, co_topo = list(coarse.carrier), list(coarse.topology)
idmap = {0: 0, 1: 1}
print("id: discrete -> sierpinski:", is_continuous_finite_map(fine_pts, fine_topo, co_pts, co_topo, idmap))
print("id: sierpinski -> discrete:", is_continuous_finite_map(co_pts, co_topo, fine_pts, fine_topo, idmap))
\end{lstlisting}

\begin{verbatim}
id: discrete -> sierpinski: True
id: sierpinski -> discrete: False
\end{verbatim}

Ayrıktaki açık $\{0\}$'ın geri çekimi yine $\{0\}$'dır; bu Sierpiński
$\tau = \{\emptyset, \{1\}, \{0,1\}\}$ içinde açık olmadığından ``kaba $\to$
ince'' yönü süreklilik kaybeder.

%-------------------------------------------------------------------
\section{Alıştırmalar}

\subsection*{Kodlama}

\begin{enumerate}[label=K\arabic*.]
\item Özel bir topolojide $f(0)=0, f(1)=0, f(2)=1$ fonksiyonunun sürekli
      olup olmadığını kontrol edin.
\item Sierpiński uzayından kendisine giden dört fonksiyonu test edin.
\item \texttt{finite\_homeomorphism\_result} ile iki ayrık uzayın homeomorf olduğunu doğrulayın.
\item \texttt{make\_topology} ile $\{1,2,3\}$ üzerinde $\{\alpha\}$, $\{\alpha,\beta\}$
      açıklarıyla zincir topolojileri kurun; zinciri koruyan $f$ ve $g$ tanımlayıp
      $f$, $g$ ve $g\circ f$'in üçünün de sürekli olduğunu
      \texttt{is\_continuous\_finite\_map} ile doğrulayın. \ipucu{Teorem~2.2'nin sayısal kontrolü.}
\end{enumerate}

\subsection*{Teori}

\begin{enumerate}[label=T\arabic*.]
\item Her sabit fonksiyonun sürekli olduğunu ispatlayın.
\item $f \circ g$ bileşkesinin sürekli olduğunu ispatlayın.
\item $f: X \to Y$ sürekli ve $X$ kompakt ise $f(X)$'in kompakt olduğunu
      ispatlayın. \ipucu{$f(X)$'in açık örtüsünü $f$ ile geri çekin, $X$'in
      kompaktlığıyla sonlu alt-örtü alın, görüntüye geri itin.}
\end{enumerate}
